\documentclass[11pt,letterpaper]{article}
\usepackage[utf8]{inputenc}
\usepackage{geometry}
\usepackage{amsmath}
\usepackage{setspace}
\usepackage{url}
\usepackage{enumitem}
\usepackage{hyperref}

\geometry{left=1.0in,right=1.0in,top=1.0in,bottom=1.0in}
\onehalfspacing

\setlength{\parindent}{0pt}
\setlength{\parskip}{1.0em plus 0.2em minus 0.2em}

\begin{document}

\thispagestyle{empty}

\begin{center}
    \textbf{\large IN THE UNITED STATES PATENT AND TRADEMARK OFFICE}
\end{center}

\vspace{0.3cm}
\noindent\begin{tabular}{p{2.3in}p{4.0in}}
\textbf{Applicant / Inventor:} & Reinaldo Javier Aguiar Marcano \\
\textbf{Application Number:} & 19/561,223 \\
\textbf{Filing Date:} & March 9, 2026 \\
\textbf{Confirmation Number:} & 1850 \\
\textbf{Art Unit / Examiner:} & TC 2100 / ANDREA BENNETT \\
\textbf{Office of Petitions:} & Alesia M. Brown, Attorney Advisor \\
\textbf{Customer Number:} & 226458 \\
\textbf{Entity Status:} & SMALL ENTITY (37 CFR \S{} 1.27) \\
\textbf{Title of Invention:} & SYSTEM AND METHOD FOR N-GRAM PARALLEL DATA ENCODING \\
\end{tabular}

\vspace{0.4cm}
\begin{center}
    \textbf{\large APPLICANT RESPONSE TO NOTICE OF INCOMPLETE REPLY \\
    AND NOTICE TO FILE MISSING PARTS OF NONPROVISIONAL APPLICATION} \\
    \textbf{\normalsize (SUBMISSION OF EXTENSION OF TIME PETITION, CONDITIONAL SUSPENSION, \\
    AND RESPECTFUL REQUEST FOR DEPOSIT ALLOCATION UNDER 37 CFR \S{} 1.25(b))}
\end{center}

\noindent\textbf{Mail Stop Missing Parts} \\
Commissioner for Patents \\
P.O. Box 1450 \\
Alexandria, VA 22313-1450

\vspace{0.2cm}
Sir:

In responsive reply to the \textit{Notice of Incomplete Reply} mailed August 5, 2026, the \textit{Notice to File Missing Parts} mailed April 7, 2026, and the \textit{Decision on Petition} (PETDEC) mailed October 1, 2026, the pro se Applicant, Reinaldo Javier Aguiar Marcano, respectfully submits this comprehensive formalities response and fee reconciliation.

\section*{I. STATEMENT OF REAFFIRMED ENTITY STATUS (37 CFR \S{} 1.27)}

Applicant hereby formally reaffirms entitlement to \textbf{Small Entity status} under 37 CFR \S{} 1.27 and 35 U.S.C. \S{} 123. Applicant's verified Statement of Small Entity Status was filed on August 1, 2026, and officially recognized by the Office in the August 5, 2026 Notice of Incomplete Reply and Fee Determination Record. All statutory fees are governed by the Small Entity fee schedule.

\section*{II. PROCEDURAL POSTURE \& CURED FORMALITIES STATUS}

Applicant respectfully notes the following procedural milestones established in the Image File Wrapper:
\begin{enumerate}[leftmargin=*]
    \item \textbf{Cured Technical Informalities:} In response to the April 7, 2026 Notice to File Missing Parts, Applicant timely submitted replacement Claims 1--6 commencing on a separate sheet and a replacement Abstract commencing on a separate sheet on July 31, 2026. The Office officially confirmed entry of these compliant technical papers in the August 5, 2026 Notice of Incomplete Reply.
    \item \textbf{Inventor's Declaration:} An executed Inventor's Oath or Declaration under 37 CFR \S{} 1.63 was filed on July 24, 2026, and acknowledged in the record.
    \item \textbf{Outstanding Pre-Exam Item:} The August 5, 2026 Notice of Incomplete Reply establishes that the sole remaining items required to advance the application are the pre-examination statutory fee balance and the requisite extension of time under 37 CFR \S{} 1.136(a).
\end{enumerate}

\section*{III. STATUTORY FEE RECONCILIATION, REQUEST FOR DEPOSIT ALLOCATION (37 CFR \S{} 1.25(b)), AND MACRO PORTFOLIO EQUITY}

\subsection*{A. Macro Portfolio Fee Equity Context}
Applicant respectfully directs the Office of Patent Application Processing and the Office of Finance to the overarching financial realities of this prosecution. Across the 102 nonprovisional patent applications comprising Applicant's master portfolio, the USPTO Office of Finance is currently holding approximately \$35,000.00 to \$42,000.00 USD (with \$32,758.00 USD already empirically verified and credited across 79 examined notices in the Image File Wrappers) in unapplied funds on deposit, without interest and without refund.

Viewing this individual application through an isolated, single-docket lens distorts the financial record by suggesting that Applicant has failed to remit fees. In truth, Applicant has already remitted substantial capital that remains on deposit with the agency. Fundamental administrative equity and MPEP \S{} 509 require that unapplied deposit credits be properly accounted for and applied to maintain procedural timeliness, rather than assessing progressive extension fees that compound to over \$400,000 across the portfolio without prior fee reconciliation.

\subsection*{B. Respectful Request for Allocation of Unapplied Deposit Funds (37 CFR \S{} 1.25(b) \& MPEP \S{} 509)}
Official Office records (including the August 5, 2026 Notice of Incomplete Reply and Fee Worksheet WFEE) confirm that the Office of Finance currently holds \textbf{\$416.00 in unapplied funds on deposit} in this application.

Pursuant to 37 CFR \S{} 1.25(b), MPEP \S{} 509, and MPEP \S{} 710.02(e), Applicant hereby respectfully requests that the Office allocate said unapplied funds on deposit in the following order of procedural priority:
\begin{enumerate}[leftmargin=*]
    \item \textbf{PRIORITY 1 (Timeliness Preservation):} The gross fee for a Five-Month Extension of Time under 37 CFR \S{} 1.136(a) is \$1,358.00 (Small Entity, Fee Code 2255). Pursuant to MPEP \S{} 710.02(e), the \$276.00 extension fee previously paid on July 31, 2026 under Fee Code 2252 must be credited against the total fee due, leaving a net incremental extension fee of \textbf{\$1,082.00}. Applicant respectfully requests that the Office apply the entire \textbf{\$416.00 unapplied deposit credit} directly toward this incremental extension fee as Priority 1, reducing the net extension shortfall to \textbf{\$666.00}. Crediting this amount to the extension fee first establishes that this response is timely submitted within the statutory reply window ending November 9, 2026 (35 U.S.C. \S{} 21 and 37 CFR \S{} 1.7(a));
    \item \textbf{PRIORITY 2 (Suspension Fee Satisfaction):} Satisfy the \textbf{\$60.00 processing fee} under 37 CFR \S{} 1.17(i) (Fee Code 2830) for the concurrent Request for Suspension under 37 CFR \S{} 1.103(a);
    \item \textbf{PRIORITY 3--5 (Pre-Exam Base Fee Deficiency Deferred under 37 CFR \S{} 1.53(f)):} The remaining pre-examination base fee deficiency of \textbf{\$486.00} (comprising the \$70.00 basic filing fee, \$308.00 search fee, \$352.00 examination fee, and \$172.00 non-DOCX surcharge, less the \$416.00 credit) is not due to be resolved immediately to preserve timeliness. Because this reply is timely filed within the 5-month extension period, Applicant respectfully requests that the Office issue a formal Notice of Fee Deficiency under 37 CFR \S{} 1.53(f) and MPEP \S{} 506 for any remaining base fee balance, or hold collection in abeyance pursuant to the concurrent Request for Suspension under 37 CFR \S{} 1.103(a).
\end{enumerate}

\section*{IV. CONCLUSION}

All formality requirements having been satisfied, Applicant respectfully requests that the Office of Patent Application Processing (OPAP) advance the application to Technology Center 2100 for examination.

\vspace{0.6cm}
\begin{singlespace}
\noindent Respectfully submitted, \par\vspace{0.8cm}
\underline{\textbf{/Reinaldo Javier Aguiar Marcano/}} \\
\textbf{Reinaldo Javier Aguiar Marcano} \\
Inventor and Pro Se Applicant \\
Mailing Address: 2302 Britton Ridge Drive, Katy, Texas 77494 \\
Telephone: (202) 599-3979 \\
Email: \texttt{rj.aguiar@gmail.com} (USPTO Account) \\
Email: \texttt{rj.aguiar.m@gmail.com} \\
Email: \texttt{reinaldo.aguiar.marcano@isbn9781971481296.com} \\
Date: October 10, 2026
\end{singlespace}

\end{document}
